\section{Bias and Ethics}
\label{sec:ethics}

The results carry one bias that is the subject of the study and several that are not. Textbook diagrams are drawn to look right, so a model can learn to read marks and be scored as if it read geometry. Grading on the drawing removes that bias from the score but hands the judgement to the tolerances; we treated them as an instrument, calibrated, gated against a human reading and frozen. In a figure not drawn to scale the marks, not the drawing, tell the truth; the \nNtsTestDiagrams{} flagged test diagrams score \fNtsNotToScaleId{} and \fNtsNotToScaleHo{} against \fNtsToScaleId{} and \fNtsToScaleHo{} for the rest, too few for a claim, so they are reported as their own category rather than left to lower a score silently. The data is from English-language textbooks and one person labelled the gold set; our conclusions stop at that domain.

The uses that motivate the work act on the model's answer. A grader that misreads a figure penalises a correct observation; a tutor explains the wrong problem; a screen reader describing a diagram to a blind student misinforms someone who cannot check the page. Three safeguards answer this. The model returns a calibrated probability per relation (expected calibration error \cite{guo2017calibration} \eceFtHo{} on the unseen source), so a downstream system can abstain. The verifier is deterministic: given coordinates, any prediction can be re-measured instead of believed, and the stored margin tells how near the case sits to the line. And what comes out is a relation rather than a grade, so a person stays in the loop.

What the verifier does not protect against must be said plainly: it needs annotated coordinates, so in deployment it audits a model rather than guaranteeing an answer; a figure not drawn to scale defeats it by design; and a model rewarded for calling the aligned edges of two unrelated figures parallel is rewarded for a reading no teacher would accept.
